%%%%%%%%%%%%%%%%%%%%%%%%%%%%%%%%%%%%%%%%%%%%%%%%%%%%%%%%%%%%%%%%%%%%%%%%%%%%%%%
%                       CARREGA DE LA CLASSE DE DOCUMENT                      %
%                                                                             %
% Les opcions admissibles son:                                                %
%      12pt / 11pt            (cos dels tipus de lletra; no feu servir 10pt)  %
%                                                                             %
% catalan/spanish/english     (llengua principal del treball)                 %
%                                                                             % 
% french/italian/german...    (si necessiteu fer servir alguna altra llengua) %
%                                                                             %
% listoffigures               (El document inclou un Index de figures)        %
% listoftables                (El document inclou un Index de taules)         %
% listofquadres               (El document inclou un Index de quadres)        %
% listofalgorithms            (El document inclou un Index d'algorismes)      %
%                                                                             %
%%%%%%%%%%%%%%%%%%%%%%%%%%%%%%%%%%%%%%%%%%%%%%%%%%%%%%%%%%%%%%%%%%%%%%%%%%%%%%%

\documentclass[11pt,spanish,listoffigures,listoftables]{tfgetsinf}

%%%%%%%%%%%%%%%%%%%%%%%%%%%%%%%%%%%%%%%%%%%%%%%%%%%%%%%%%%%%%%%%%%%%%%%%%%%%%%%
%                     CODIFICACIO DEL FITXER FONT                             %
%                                                                             %
%    windows fa servir normalment 'ansinew'                                   %
%    amb linux es possible que siga 'latin1' o 'latin9'                       %
%    Pero el mes recomanable es fer servir utf8 (unicode 8)                   %
%                                          (si el vostre editor ho permet)    % 
%%%%%%%%%%%%%%%%%%%%%%%%%%%%%%%%%%%%%%%%%%%%%%%%%%%%%%%%%%%%%%%%%%%%%%%%%%%%%%%

\usepackage[utf8]{inputenc} 
%%%%%%%%%%%%%%%%%%%%%%%%%%%%%%%%%%%%%%%%%%%%%%%%%%%%%%%%%%%%%%%%%%%%%
% Para conseguir que la tabla de contenido no salga en rojo
%%%%%%%%%%%%%%%%%%%%%%%%%%%%%%%%%%%%%%%%%%%%%%%%%%%%%%%%%%%%%%%%%%%%%

\hypersetup{ colorlinks=true, linkcolor=black, urlcolor=cyan, }

%%%%%%%%%%%%%%%%%%%%%%%%%%%%%%%%%%%%%%%%%%%%%%%%%%%%%%%%%%%%%%%%%%%%%%%%%%%%%%%
%                        ALTRES PAQUETS I DEFINICIONS                         %
%                                                                             %
% Carregueu aci els paquets que necessiteu i declareu les comandes i entorns  %
%                                          (aquesta seccio pot ser buida)     %
%%%%%%%%%%%%%%%%%%%%%%%%%%%%%%%%%%%%%%%%%%%%%%%%%%%%%%%%%%%%%%%%%%%%%%%%%%%%%%%



%%%%%%%%%%%%%%%%%%%%%%%%%%%%%%%%%%%%%%%%%%%%%%%%%%%%%%%%%%%%%%%%%%%%%%%%%%%%%%%
%                        DADES DEL TREBALL                                    %
%                                                                             %
% titol, alumne, tutor i curs academic                                        %
%%%%%%%%%%%%%%%%%%%%%%%%%%%%%%%%%%%%%%%%%%%%%%%%%%%%%%%%%%%%%%%%%%%%%%%%%%%%%%%

\title{Transferencia de modelos MEG de contexto largo a EEG para clasificación de palabras en lectura natural
}
\author{Ricardo Díaz Peris}
\tutor{Vicent Juan Botti Navarro, Alfons Juan Císcar}
\curs{2025-2026}

%%%%%%%%%%%%%%%%%%%%%%%%%%%%%%%%%%%%%%%%%%%%%%%%%%%%%%%%%%%%%%%%%%%%%%%%%%%%%%%
%                     PARAULES CLAU/PALABRAS CLAVE/KEY WORDS                  %
%                                                                             %
% Independentment de la llengua del treball, s'hi han d'incloure              %
% les paraules clau i el resum en els tres idiomes                            %
%%%%%%%%%%%%%%%%%%%%%%%%%%%%%%%%%%%%%%%%%%%%%%%%%%%%%%%%%%%%%%%%%%%%%%%%%%%%%%%

\keywords{Descodificació neuronal; EEG; MEG; aprenentatge per transferència; lectura natural; classificació de paraules; models de context llarg; brain-to-text} % Paraules clau 
         {Decodificación neuronal; EEG; MEG; aprendizaje por transferencia; lectura natural; clasificación de palabras; modelos de contexto largo; brain-to-text}              % Palabras clave
         {Neural decoding; EEG; MEG; transfer learning; natural reading; word classification; long-context models; brain-to-text}        % Key words

%%%%%%%%%%%%%%%%%%%%%%%%%%%%%%%%%%%%%%%%%%%%%%%%%%%%%%%%%%%%%%%%%%%%%%%%%%%%%%%
%                              INICI DEL DOCUMENT                             %
%%%%%%%%%%%%%%%%%%%%%%%%%%%%%%%%%%%%%%%%%%%%%%%%%%%%%%%%%%%%%%%%%%%%%%%%%%%%%%%

\begin{document}

%%%%%%%%%%%%%%%%%%%%%%%%%%%%%%%%%%%%%%%%%%%%%%%%%%%%%%%%%%%%%%%%%%%%%%%%%%%%%%%
%              RESUMS DEL TFG EN VALENCIA, CASTELLA I ANGLES                  %
%%%%%%%%%%%%%%%%%%%%%%%%%%%%%%%%%%%%%%%%%%%%%%%%%%%%%%%%%%%%%%%%%%%%%%%%%%%%%%%

\begin{abstract}
La decodificación de palabras a partir de señales cerebrales no invasivas presenta importantes dificultades debidas al bajo nivel de señal útil, la variabilidad entre sujetos y las diferencias entre modalidades como MEG y EEG. En este trabajo se adapta el modelo MEG-XL, diseñado originalmente para el aprendizaje de representaciones de contexto largo en MEG, a un escenario de EEG durante lectura natural. La adaptación se centra en construir un pipeline reproducible para evaluar si las representaciones aprendidas en MEG pueden transferirse a EEG en tareas de clasificación de palabras.
Para ello, se integra la arquitectura de MEG-XL con un nuevo adaptador diferentes datasets de EEG, que transforma los registros de este tipo de sensores y las anotaciones de lectura en segmentos alineados a palabras. El pipeline incorpora re-muestreo, normalización, extracción de ventanas temporales por palabra, máscaras de sensores y codificación explícita del tipo de sensor EEG. A partir de estos segmentos, el modelo obtiene representaciones neuronales contextualizadas mediante un Transformer con atención criss-cross y las compara con embeddings lingüísticos generados por T5 mediante una pérdida contrastiva.
El objetivo principal es analizar la viabilidad técnica y experimental de trasladar modelos de contexto largo entrenados en MEG a datos EEG para comenzar a investigar una solución cerrada de brain-to-text que incluya EEG. Este enfoque permite estudiar las barreras prácticas de la transferencia entre modalidades, identificar los puntos críticos del preprocesamiento y sentar una base para futuros modelos multimodales M/EEG aplicados a decodificación de lenguaje natural.
\end{abstract}
\begin{abstract}[spanish]
La descodificació de paraules a partir de senyals cerebrals no invasives presenta importants dificultats degudes al baix nivell de senyal útil, la variabilitat entre subjectes i les diferències entre modalitats com MEG i EEG. En este treball s’adapta el model MEG-XL, dissenyat originalment per a l'aprenentatge de representacions de context llarg en MEG, a un escenari d'EEG durant lectura natural. L'adaptació se centra a construir un pipeline reproduïble per a avaluar si les representacions apreses en MEG poden transferir-se a EEG en tasques de classificació de paraules.
Per això, s'integra l'arquitectura de MEG-XL amb un nou adaptador per a diferents conjunts de dades d’EEG, que transforma els registres d'aquest tipus de sensors i les anotacions de lectura en segments alineats amb paraules. El pipeline incorpora remostratge, normalització, extracció de finestres temporals per paraula, màscares de sensors i codificació explícita del tipus de sensor EEG. A partir d'aquestos segments, el model obté representacions neuronals contextualitzades mitjançant un Transformer amb atenció criss-cross i les compara amb embeddings lingüístics generats per T5 mitjançant una pèrdua contrastiva.
L'objectiu principal és analitzar la viabilitat tècnica i experimental de traslladar models de context llarg entrenats en MEG a dades EEG per a començar a investigar una solució tancada de brain-to-text que incloga EEG. Este enfocament permet estudiar les barreres pràctiques de la transferència entre modalitats, identificar els punts crítics del preprocessament i assentar una base per a futurs models multimodals M/EEG aplicats a la descodificació de llenguatge natural.
\end{abstract}
\begin{abstract}[english]
Decoding words from non-invasive brain signals presents significant challenges due to the low level of useful signal, inter-subject variability, and differences between modalities such as MEG and EEG. In this work, the MEG-XL model—originally designed for learning long-range context representations in MEG—is adapted to an EEG setting during natural reading. The adaptation focuses on building a reproducible pipeline to evaluate whether representations learned in MEG can be transferred to EEG in word classification tasks.
To this end, the MEG-XL architecture is integrated with a new adapter for various EEG datasets, which transforms EEG recordings and reading annotations into word-aligned segments. The pipeline incorporates resampling, normalization, word-based temporal window extraction, sensor masking, and explicit encoding of the EEG sensor type. From these segments, the model obtains contextualized neural representations using a Transformer with cross-attention and compares them with linguistic embeddings generated by T5 via a contrastive loss.
The main objective is to analyze the technical and experimental feasibility of transferring long-context models trained on MEG data to EEG data in order to begin investigating an end-to-end brain-to-text solution that incorporates EEG. This approach allows us to study the practical barriers to cross-modal transfer, identify critical points in preprocessing, and lay the groundwork for future multimodal M/EEG models applied to natural language decoding.
\end{abstract}

%%%%%%%%%%%%%%%%%%%%%%%%%%%%%%%%%%%%%%%%%%%%%%%%%%%%%%%%%%%%%%%%%%%%%%%%%%%%%%%
%                              CONTINGUT DEL TREBALL                          %
%%%%%%%%%%%%%%%%%%%%%%%%%%%%%%%%%%%%%%%%%%%%%%%%%%%%%%%%%%%%%%%%%%%%%%%%%%%%%%%

\mainmatter

%%%%%%%%%%%%%%%%%%%%%%%%%%%%%%%%%%%%%%%%%%%%%%%%%%%%%%%%%%%%%%%%%%%%%%%%%%%%%%%
%                                  INTRODUCCIO                                %
%%%%%%%%%%%%%%%%%%%%%%%%%%%%%%%%%%%%%%%%%%%%%%%%%%%%%%%%%%%%%%%%%%%%%%%%%%%%%%%

\chapter{Introducción}

????? ????????????? ????????????? ????????????? ????????????? ????????????? 


\section{Motivación}

????? ????????????? ????????????? ????????????? ????????????? ????????????? 

\section{Objetivos}

????? ????????????? ????????????? ????????????? ????????????? ?????????????

\section{Estructura de la memoria}

????? ????????????? ????????????? ????????????? ????????????? ????????????? 

%\section{Notes bibliografiques} %%%%% Opcional

%????? ????????????? ????????????? ????????????? ????????????? ?????????????

%%%%%%%%%%%%%%%%%%%%%%%%%%%%%%%%%%%%%%%%%%%%%%%%%%%%%%%%%%%%%%%%%%%%%%%%%%%%%%%
%                         CAPITOLS (tants com calga)                          %
%%%%%%%%%%%%%%%%%%%%%%%%%%%%%%%%%%%%%%%%%%%%%%%%%%%%%%%%%%%%%%%%%%%%%%%%%%%%%%%

% Fragmento preparado para integrarse directamente en la memoria principal.
% Paquetes recomendados en el preámbulo: amsmath, amssymb, array.
% Las claves bibliográficas utilizadas se incluyen en referencias_marco_teorico.bib.

\chapter{Marco teórico}

\section{Interfaces cerebro-computador}

% Una interfaz cerebro-computador, conocida habitualmente como BCI por sus siglas en inglés \textit{Brain-Computer Interface}, es un sistema que establece un canal de comunicación entre la actividad cerebral de una persona y un dispositivo externo \cite{wolpaw2002bci}. A diferencia de las interfaces convencionales, como un teclado, un ratón o una pantalla táctil, una BCI intenta interpretar directamente determinados patrones de actividad neuronal y transformarlos en una salida útil. Esta salida puede consistir, por ejemplo, en mover un cursor, seleccionar una letra, controlar una prótesis o generar una secuencia de texto.

Las interfaces cerebro-computador (BCIs) son sistemas que establecen un canal de comunicación entre la actividad cerebral de una persona y un dispositivo externo 

% El interés por este tipo de sistemas ha aumentado considerablemente durante los últimos años. Además del trabajo desarrollado en universidades y centros clínicos, la aparición de iniciativas privadas de gran visibilidad, como Neuralink, ha contribuido a acercar las interfaces cerebro-computador al debate público \cite{musk2019neuralink}. Sin embargo, más allá de las expectativas generadas alrededor de estas tecnologías, su desarrollo sigue planteando dificultades importantes relacionadas con la adquisición de señales, la seguridad de los participantes, la cantidad de datos disponibles y la capacidad de los modelos para generalizar entre personas y sesiones.

% De forma general, una BCI se puede dividir en varias etapas. En primer lugar, se registra la actividad cerebral mediante una modalidad de adquisición, como electroencefalografía, magnetoencefalografía o electrocorticografía. Después, la señal obtenida se preprocesa para reducir ruido, artefactos y diferencias de escala. Sobre estos datos se extrae una representación que conserve la información relevante para la tarea. Finalmente, un modelo de decodificación transforma dicha representación en una predicción y, en los sistemas en línea, esta predicción se utiliza para actuar sobre un dispositivo o proporcionar realimentación al usuario.

% No todas las BCI persiguen el mismo objetivo. Algunas se basan en respuestas provocadas por estímulos externos, como los potenciales P300 o los potenciales visuales evocados de estado estable. Otras intentan reconocer actividad generada voluntariamente por el usuario, como la imaginación motora o el habla imaginada. También existen sistemas pasivos que no buscan ejecutar una orden, sino estimar estados como la carga cognitiva, la atención o la fatiga. Dentro de esta variedad de aplicaciones, la decodificación cerebral del lenguaje es una de las líneas de investigación más ambiciosas, ya que pretende recuperar información lingüística directamente a partir de la actividad neuronal.

\subsection{Decodificación cerebral del lenguaje}

% La decodificación cerebral del lenguaje consiste en inferir información relacionada con el habla o la lectura a partir de señales cerebrales. Dependiendo del experimento, esta información puede corresponder a fonemas, sílabas, palabras, frases completas o propiedades semánticas del contenido procesado por el participante. Cuando el objetivo final es reconstruir texto de forma continua, se suele emplear el término \textit{brain-to-text}. No obstante, antes de alcanzar una transcripción abierta y completa, gran parte de la investigación se centra en tareas intermedias más controladas, como la detección de presencia de habla, la clasificación de fonemas o la identificación de palabras dentro de un vocabulario limitado \cite{ozdogan2025libribrain,jayalath2025unlocking}.

% Es importante distinguir entre las diferentes situaciones en las que el lenguaje puede estar presente. En el habla producida o \textit{overt speech}, la persona pronuncia realmente una palabra o una frase. Esta condición permite comprobar fácilmente el contenido generado, pero introduce artefactos musculares relacionados con el movimiento de la boca, la lengua y la mandíbula. En el habla imaginada, también denominada \textit{covert speech}, \textit{inner speech} o habla interna, la persona imagina que pronuncia el contenido sin mover voluntariamente los articuladores ni generar sonido. Esta segunda situación resulta especialmente interesante para personas que no pueden hablar, aunque también es más difícil de registrar y verificar \cite{panachakel2021covert,alzahrani2024systematic}.

% Una tercera posibilidad es estudiar la percepción del habla. En estos experimentos, el participante escucha palabras, frases o narraciones mientras se registra su actividad cerebral. Al existir un estímulo acústico conocido, es posible alinear con precisión los eventos lingüísticos con la señal neuronal. Esta característica ha permitido crear conjuntos de datos de gran tamaño basados en audiolibros y habla continua, como LibriBrain, en el que un participante escucha más de cincuenta horas de narraciones en inglés mientras se registra MEG \cite{ozdogan2025libribrain}. Aunque escuchar habla no equivale a producirla o imaginarla, constituye un escenario adecuado para estudiar representaciones lingüísticas y desarrollar modelos de decodificación no invasiva a gran escala.

% Por último, el lenguaje también puede estudiarse durante la lectura. En este caso, la actividad cerebral se registra mientras el participante procesa palabras presentadas de forma aislada o dentro de frases. La lectura natural introduce una dificultad adicional: cada persona fija la mirada sobre las palabras durante intervalos diferentes y puede omitir o volver a leer fragmentos del texto. Por este motivo, algunos conjuntos de datos combinan EEG y seguimiento ocular. El \textit{eye-tracking} permite localizar el inicio y la duración de cada fijación y, a partir de esta información, extraer segmentos de EEG asociados a palabras concretas. ZuCo es uno de los ejemplos más conocidos de este tipo de adquisición, ya que proporciona EEG de alta densidad y seguimiento ocular simultáneo durante la lectura natural de frases \cite{hollenstein2020zuco2}.

% La señal relacionada con una palabra no se limita necesariamente al instante exacto en el que esta aparece. El procesamiento lingüístico presenta dependencias a diferentes escalas temporales. Los fonemas se agrupan en sílabas, las sílabas forman palabras y las palabras adquieren significado dentro de frases y discursos. Además, la respuesta cerebral a una palabra puede depender de las palabras anteriores, de las expectativas del lector u oyente y del estado general del participante. Por ello, los enfoques recientes no solo analizan ventanas aisladas, sino secuencias de segmentos neuronales consecutivos. Esta incorporación de contexto permite al modelo utilizar relaciones entre palabras y patrones cerebrales que no son visibles cuando cada segmento se procesa de manera independiente \cite{dascoli2024words,jayalath2026megxl}.

% La decodificación del lenguaje no invasiva continúa siendo un problema especialmente difícil. Las señales registradas fuera del cráneo tienen una relación señal-ruido reducida, los conjuntos de datos suelen contener pocas horas por participante y existen grandes diferencias entre sujetos, dispositivos y sesiones. A esto se añade el tamaño del espacio de salida: distinguir entre dos estados es mucho más sencillo que elegir una palabra entre cientos o generar una frase completa. Como consecuencia, la clasificación de palabras dentro de vocabularios cerrados sigue siendo una tarea fundamental para medir el progreso de los modelos antes de abordar sistemas \textit{brain-to-text} completamente abiertos.

\subsection{Métodos invasivos y no invasivos}

% Las técnicas utilizadas para registrar actividad cerebral se pueden dividir, de forma general, en métodos invasivos y no invasivos. Los métodos invasivos requieren una intervención quirúrgica para situar los sensores sobre la superficie del cerebro o dentro del tejido cerebral. Entre ellos se encuentran la electrocorticografía, conocida como ECoG, y los conjuntos de microelectrodos intracorticales.

% La principal ventaja de estos métodos es la calidad de la señal. Al encontrarse cerca de las poblaciones neuronales que generan la actividad, los electrodos registran señales con mayor amplitud, mejor resolución espacial y menos atenuación que los sensores externos. Esta proximidad ha permitido obtener resultados muy destacados en decodificación motora y del habla. En particular, algunos sistemas invasivos han conseguido reconstruir texto o sintetizar voz en personas con parálisis a partir de actividad cortical relacionada con los intentos de habla \cite{jayalath2025unlocking}. Sin embargo, estas prestaciones se obtienen a cambio de riesgos quirúrgicos, posibles infecciones, degradación de los electrodos y una cobertura limitada a las regiones implantadas. Por este motivo, su utilización suele restringirse a casos clínicos concretos y entornos de investigación controlados.

% Los métodos no invasivos registran la actividad cerebral sin necesidad de cirugía. Dentro de este grupo se encuentran EEG, MEG, resonancia magnética funcional y espectroscopia funcional del infrarrojo cercano. Estas técnicas reducen considerablemente los riesgos para el participante y permiten realizar estudios repetidos con mayor facilidad. Su principal limitación es que las señales deben atravesar diferentes tejidos o medirse a cierta distancia de su origen, lo que disminuye su resolución y su relación señal-ruido.

% En la decodificación del lenguaje, EEG y MEG destacan por su elevada resolución temporal. Ambas modalidades permiten seguir cambios cerebrales del orden de milisegundos, una propiedad importante para estudiar procesos tan rápidos como la percepción de fonemas o palabras. La resonancia magnética funcional ofrece una mejor resolución espacial, pero su respuesta hemodinámica es mucho más lenta y dificulta la separación temporal de palabras consecutivas. EEG y MEG constituyen, por tanto, un compromiso especialmente interesante entre seguridad y capacidad para representar la dinámica temporal del lenguaje.

% La distinción entre métodos invasivos y no invasivos también afecta al tipo de modelos que pueden entrenarse. Los sistemas invasivos disponen de señales más localizadas y, en algunos estudios, de muchas horas registradas para un mismo paciente. En cambio, los modelos no invasivos deben aprender a trabajar con señales más débiles, sensores distribuidos por toda la cabeza y una variabilidad elevada. En este contexto, el preentrenamiento, la transferencia de aprendizaje y la utilización de contexto temporal aparecen como estrategias relevantes para compensar la menor calidad de la señal.

\section{Electroencefalografía y magnetoencefalografía}

% La electroencefalografía y la magnetoencefalografía son dos técnicas no invasivas que miden manifestaciones diferentes de una misma actividad neuronal. Las corrientes postsinápticas generadas por poblaciones de neuronas producen tanto potenciales eléctricos como campos magnéticos. EEG registra las diferencias de potencial que llegan al cuero cabelludo, mientras que MEG mide los campos magnéticos producidos alrededor de la cabeza. Aunque ambas modalidades ofrecen una resolución temporal elevada, sus propiedades físicas, su distribución espacial y sus sistemas de adquisición son diferentes.

% Esta diferencia es especialmente importante cuando se pretende transferir un modelo entrenado con una modalidad a la otra. Un modelo no recibe directamente la actividad de las fuentes cerebrales, sino una mezcla de dichas fuentes observada a través de una disposición concreta de sensores. Por tanto, cambiar de MEG a EEG no consiste únicamente en modificar el número de canales. También cambian la amplitud, la geometría, el tipo de ruido, la referencia y la relación entre cada sensor y las fuentes neuronales.

\subsection{Electroencefalografía}

% La electroencefalografía, conocida como EEG, registra diferencias de potencial eléctrico mediante electrodos colocados sobre el cuero cabelludo. Estas diferencias se producen principalmente por la actividad sincronizada de grandes poblaciones de neuronas corticales. La amplitud registrada suele encontrarse en el orden de los microvoltios, por lo que el sistema de adquisición necesita amplificar señales muy pequeñas y mantener un buen contacto entre los electrodos y la piel.

% La colocación de los electrodos suele seguir sistemas estandarizados, como el sistema internacional 10--20 y sus extensiones 10--10 o 10--5. Estos sistemas asignan nombres a las posiciones en función de la región aproximada de la cabeza: frontal, central, temporal, parietal y occipital. En montajes de alta densidad se pueden utilizar 64, 128 o incluso más canales. Un mayor número de electrodos proporciona una observación espacial más detallada, aunque también incrementa el tiempo de preparación, la posibilidad de canales defectuosos y el coste computacional.

% EEG es una medida relativa. Cada canal representa la diferencia entre un electrodo y una referencia, o una combinación de referencias. La elección de esta referencia afecta a la forma de la señal y debe tenerse en cuenta al comparar conjuntos de datos. Entre las opciones habituales se encuentran la referencia a un electrodo concreto, la referencia promedio de todos los canales o configuraciones bipolares. En conjuntos de datos recogidos con sistemas diferentes, una estrategia de referencia homogénea puede reducir parte de las diferencias, pero no elimina por completo la variabilidad entre montajes.

% Una de las principales ventajas de EEG es su resolución temporal. Las señales se registran normalmente a frecuencias de cientos o miles de muestras por segundo, lo que permite estudiar respuestas rápidas asociadas a estímulos visuales, auditivos o lingüísticos. También es una tecnología relativamente económica, portátil y compatible con experimentos naturales. Estas características explican su uso frecuente en BCI y en estudios de lectura, donde puede combinarse con seguimiento ocular para obtener segmentos asociados a cada fijación \cite{hollenstein2020zuco2}.

% Su principal limitación es la resolución espacial. Los potenciales eléctricos se propagan a través del cerebro, el líquido cefalorraquídeo, el cráneo y el cuero cabelludo. La conductividad de estos tejidos produce un fenómeno de difusión conocido como conducción de volumen, por el cual una misma fuente puede afectar a numerosos electrodos. Como consecuencia, la actividad observada en un canal no puede asociarse directamente a una única región cerebral.

% EEG también es sensible a diferentes tipos de artefactos. Los parpadeos y movimientos oculares generan señales de gran amplitud en los electrodos frontales. La tensión muscular, el movimiento de la mandíbula, el latido cardíaco, el desplazamiento de los cables y el contacto deficiente de los electrodos pueden ocultar la actividad neuronal. En tareas de habla producida, los artefactos musculares son especialmente relevantes. En lectura natural, los movimientos oculares forman parte de la propia tarea, de modo que su tratamiento debe realizarse con cuidado para no eliminar información relacionada indirectamente con el procesamiento lingüístico.

% La actividad EEG también suele describirse mediante bandas de frecuencia. De manera aproximada, se distinguen las bandas delta, theta, alpha, beta y gamma. Estas bandas no representan componentes completamente independientes, pero sirven para resumir cambios oscilatorios relacionados con diferentes procesos. En estudios de habla imaginada se han utilizado combinaciones muy distintas de frecuencias, desde bandas bajas hasta gamma, sin que exista un intervalo universalmente óptimo \cite{panachakel2021covert,alzahrani2024systematic}. Por este motivo, la frecuencia de muestreo y el filtrado deben considerarse hiperparámetros relevantes cuando se adaptan modelos a nuevos conjuntos de EEG.

\subsection{Magnetoencefalografía}

% La magnetoencefalografía, conocida como MEG, registra los campos magnéticos generados por la actividad neuronal. Estos campos son extremadamente débiles, normalmente del orden de femtoteslas, por lo que se necesitan sensores muy sensibles y salas protegidas frente a interferencias magnéticas externas. Los sistemas tradicionales utilizan dispositivos superconductores de interferencia cuántica, conocidos como SQUID. También se están desarrollando sistemas basados en magnetómetros de bombeo óptico, que permiten situar los sensores más cerca de la cabeza y ofrecen nuevas posibilidades para dispositivos más flexibles.

% Los sistemas MEG suelen incluir magnetómetros y gradiometros. Los magnetómetros miden directamente la intensidad del campo magnético, mientras que los gradiometros registran diferencias espaciales del campo y ayudan a reducir interferencias externas. Un equipo MEG de cuerpo completo puede contener más de trescientos canales distribuidos alrededor de la cabeza. Por ejemplo, LibriBrain fue registrado con un sistema MEGIN Triux Neo de 306 canales \cite{ozdogan2025libribrain}.

% A diferencia de los potenciales eléctricos, los campos magnéticos atraviesan el cráneo y el cuero cabelludo con una distorsión menor. Esto proporciona a MEG una resolución espacial superior a la de EEG en muchas condiciones, manteniendo una resolución temporal del orden de milisegundos. Estas propiedades han convertido MEG en una modalidad especialmente utilizada para estudiar percepción auditiva, lenguaje y habla continua \cite{dash2020megphrases}.

% MEG también presenta limitaciones importantes. El equipamiento tiene un coste elevado, requiere instalaciones específicas y no es fácilmente transportable. El participante debe permanecer relativamente quieto, ya que el movimiento de la cabeza modifica la relación entre las fuentes cerebrales y los sensores. Además, los datos ocupan mucho espacio y su procesamiento suele incluir pasos específicos, como corrección del movimiento de la cabeza, detección de canales defectuosos y separación de señales de origen externo mediante técnicas como \textit{Signal Space Separation} o filtrado de Maxwell.

% Otra propiedad relevante es que los sistemas MEG no comparten necesariamente la misma disposición de sensores. Distintos fabricantes utilizan cantidades, tipos y geometrías diferentes. Incluso dentro del mismo sistema, la posición de la cabeza puede cambiar entre sesiones. Para combinar datos de varios conjuntos, los modelos deben representar de algún modo la localización, la orientación y el tipo de sensor, o transformar previamente las señales a un espacio común.

% MEG ha permitido desarrollar algunos de los conjuntos de datos más relevantes para la decodificación no invasiva del habla. La mayor calidad de la señal, unida a adquisiciones prolongadas, ha facilitado tareas de detección de habla, clasificación de fonemas y clasificación contextual de palabras. Estos resultados hacen que los modelos preentrenados con MEG constituyan una base interesante para estudiar su transferencia a EEG, aunque las diferencias entre modalidades impiden asumir que dicha transferencia será directa.

\subsection{Comparación entre EEG y MEG}

% EEG y MEG comparten su carácter no invasivo y su elevada resolución temporal, pero presentan diferencias importantes en la forma en que observan la actividad cerebral. La Tabla~\ref{tab:eeg_meg_comparison} resume las propiedades más relevantes para este trabajo.

% \begin{table}[ht]
%     \centering
%     \small
%     \caption{Comparación general entre EEG y MEG.}
%     \label{tab:eeg_meg_comparison}
%     \begin{tabular}{|p{3.1cm}|p{5.2cm}|p{5.2cm}|}
%         \hline
%         \textbf{Propiedad} & \textbf{EEG} & \textbf{MEG} \\
%         \hline
%         Magnitud registrada & Diferencias de potencial eléctrico en el cuero cabelludo. & Campos magnéticos producidos por corrientes neuronales. \\
%         \hline
%         Resolución temporal & Elevada, normalmente del orden de milisegundos. & Elevada, normalmente del orden de milisegundos. \\
%         \hline
%         Resolución espacial & Limitada por la conducción de volumen y la referencia. & Generalmente superior, con menor distorsión causada por el cráneo. \\
%         \hline
%         Sensibilidad a fuentes & Sensible a fuentes radiales y tangenciales. & Principalmente sensible a fuentes tangenciales en sistemas convencionales. \\
%         \hline
%         Artefactos frecuentes & Parpadeos, movimientos oculares, actividad muscular, referencia y contacto de electrodos. & Movimiento de cabeza, interferencias magnéticas, actividad cardíaca y muscular. \\
%         \hline
%         Portabilidad y coste & Relativamente económico y portable. & Coste elevado y necesidad de instalaciones especializadas. \\
%         \hline
%         Preparación & Colocación de electrodos y control de impedancias. & Digitalización de la cabeza, localización de sensores y control de movimiento. \\
%         \hline
%         Variabilidad entre conjuntos & Número y posición de electrodos, referencia, amplificador y frecuencia de muestreo. & Geometría, tipo de sensor, posición de la cabeza y sistema de adquisición. \\
%         \hline
%     \end{tabular}
% \end{table}

% Desde el punto de vista del aprendizaje automático, ambas modalidades suelen representarse inicialmente como una matriz de canales por tiempo. Sin embargo, esta similitud de formato puede resultar engañosa. Los canales de EEG y MEG no son intercambiables, y una misma fuente cortical se proyecta de forma diferente en cada modalidad. Además, EEG suele tener menos canales y una relación señal-ruido inferior, mientras que MEG proporciona una estructura espacial más rica.

% Estas diferencias generan un problema de cambio de dominio. Un modelo entrenado con MEG puede haber aprendido patrones que dependen de la geometría de sus sensores, de la amplitud de la señal o del tipo de ruido presente en el escáner. Al aplicarlo sobre EEG, estas relaciones dejan de cumplirse. Por ello, la transferencia requiere adaptar la entrada, normalizar la señal, representar la posición de los electrodos y, en muchos casos, reajustar parte o la totalidad del modelo.

% Aun así, EEG y MEG comparten información temporal y neurofisiológica. Ambas modalidades reflejan la actividad coordinada de poblaciones neuronales y contienen respuestas relacionadas con el procesamiento de palabras. Esta base común justifica estudiar si las representaciones aprendidas mediante preentrenamiento MEG pueden aportar información útil para EEG, especialmente en escenarios con pocos datos etiquetados. El objetivo no es asumir que ambas señales son iguales, sino determinar qué partes del conocimiento aprendido son transferibles y qué componentes deben adaptarse.

\section{Procesamiento de señales cerebrales}

% Las señales EEG y MEG se registran como series temporales multicanal. Si se utilizan \(C\) sensores y se obtienen \(T\) muestras, una grabación se puede representar como una matriz
% \[
%     \mathbf{X} \in \mathbb{R}^{C \times T}.
% \]
% Esta representación contiene simultáneamente información temporal y espacial, pero también ruido, artefactos y diferencias instrumentales. Antes de entrenar un modelo es necesario decidir qué parte de la señal se conserva, cómo se normaliza y cómo se divide en muestras.

% El procesamiento no debe entenderse únicamente como una limpieza previa. Cada decisión modifica la información que el modelo puede utilizar. Un filtro demasiado restrictivo puede eliminar componentes relacionadas con la tarea, mientras que una limpieza insuficiente puede hacer que el clasificador aprenda artefactos. Del mismo modo, una segmentación incorrecta puede asociar una etiqueta lingüística a una ventana que no contiene la respuesta cerebral correspondiente.

\subsection{Preprocesamiento de EEG y MEG}

% El primer paso suele ser comprobar la calidad de la grabación. Durante esta revisión se identifican canales planos, saturados o excesivamente ruidosos. Estos canales pueden eliminarse, interpolarse a partir de sensores vecinos o marcarse para que el modelo no los utilice. También es importante verificar la sincronización entre la señal cerebral y los eventos del experimento, especialmente cuando las ventanas se alinean con el inicio de una palabra, un estímulo auditivo o una fijación ocular.

% En EEG, el preprocesamiento puede incluir una nueva referencia de los canales. La referencia promedio es una opción común en sistemas de alta densidad, aunque su idoneidad depende de la cobertura de electrodos. También se revisan las impedancias y se eliminan canales periféricos muy afectados por actividad muscular. En MEG no existe una referencia eléctrica equivalente, pero son habituales la corrección del movimiento de la cabeza y los métodos que separan las fuentes internas y externas al casco. En LibriBrain, por ejemplo, se aplica filtrado de Maxwell, eliminación de canales defectuosos, filtros para la red eléctrica y reducción de la frecuencia de muestreo \cite{ozdogan2025libribrain}.

% El filtrado temporal busca conservar un intervalo de frecuencias útil y reducir componentes no deseadas. Un filtro paso alto elimina las variaciones muy lentas y la deriva de la línea base. Un filtro paso bajo reduce el ruido de alta frecuencia y prepara la señal para un posible remuestreo. También se utilizan filtros \textit{notch} en la frecuencia de la red eléctrica, que en Europa es de 50~Hz, y en algunos de sus armónicos. La elección de los límites debe estar relacionada con la tarea y con la nueva frecuencia de muestreo. Antes de reducir la frecuencia es necesario aplicar un filtrado antialiasing que elimine componentes superiores a la frecuencia de Nyquist.

% El remuestreo reduce la cantidad de puntos temporales y, por tanto, el coste de almacenamiento y entrenamiento. Esta reducción es especialmente importante en modelos de contexto largo, donde el número de elementos de entrada crece con la duración, el número de canales y la frecuencia de muestreo. MEG-XL, por ejemplo, utiliza señales remuestreadas a 50~Hz para poder procesar ventanas de dos minutos y medio \cite{jayalath2026megxl}. Sin embargo, una frecuencia menor también limita las componentes espectrales disponibles. En EEG no existe una frecuencia óptima universal para la decodificación de palabras, por lo que su selección debe evaluarse experimentalmente.

% La eliminación de artefactos es uno de los pasos más delicados. En EEG se emplea con frecuencia el análisis de componentes independientes, conocido como ICA, para separar fuentes con patrones espaciales y temporales distintos. Posteriormente, se identifican los componentes asociados a parpadeos, movimientos oculares, actividad cardíaca o tensión muscular. Esta identificación puede realizarse manualmente o mediante métodos automáticos como MARA. ZuCo utiliza una estrategia automática basada en ICA para eliminar componentes oculares y musculares antes de sincronizar EEG y seguimiento ocular \cite{hollenstein2020zuco2}.

% A pesar de su utilidad, la eliminación de artefactos no está exenta de riesgos. En una tarea de lectura, la actividad ocular está relacionada con el momento en que se procesa cada palabra. Eliminar toda señal correlacionada con los ojos podría suprimir parte de la estructura temporal necesaria para la alineación. De manera similar, en habla producida, un modelo puede alcanzar una precisión elevada utilizando actividad muscular en lugar de actividad cerebral. Por ello, el tratamiento de artefactos debe ajustarse al objetivo del experimento y acompañarse de controles que permitan detectar atajos no deseados.

% Después del filtrado, las señales suelen normalizarse. Una normalización global puede verse dominada por diferencias entre sesiones o por valores extremos. Por este motivo, se utilizan estrategias por canal, por grabación o por segmento. Entre ellas se encuentran la estandarización mediante media y desviación típica, la normalización robusta mediante mediana y rango intercuartílico y el recorte de valores extremos. MEG-XL divide las muestras largas en subsegmentos de tres segundos, corrige la línea base, normaliza mediante estadísticos robustos y limita los valores a un intervalo fijo \cite{jayalath2026megxl}. Este tipo de transformación reduce las diferencias de escala sin depender excesivamente de valores anómalos.

% La última etapa es la segmentación. En clasificación de palabras, cada ejemplo se obtiene normalmente a partir de una ventana alineada con un evento lingüístico. En escucha, este evento puede ser el inicio de la palabra en el audio. En lectura, puede ser el momento en que comienza una fijación sobre la palabra. La ventana puede incluir señal anterior al evento para corregir la línea base y señal posterior para capturar respuestas tardías. Cuando se procesan secuencias completas, varias ventanas consecutivas se agrupan manteniendo su orden original.

% La división entre entrenamiento, validación y prueba también forma parte del procesamiento. No es suficiente con separar ventanas aleatoriamente si segmentos cercanos pertenecen a la misma grabación o al mismo estímulo. Esto puede producir fugas de información, ya que el modelo reconoce propiedades específicas de la sesión o del contenido. Siempre que sea posible, deben mantenerse sesiones, sujetos o estímulos completos dentro de una única partición. En conjuntos donde varios participantes reciben el mismo estímulo, todas las repeticiones de dicho estímulo deben asignarse a la misma partición.

\subsection{Representación temporal y espacial de las señales}

% La forma más directa de representar EEG o MEG es utilizar la matriz original de canales por tiempo. Esta opción conserva la señal con pocas suposiciones previas y permite que el modelo aprenda sus propias características. Sin embargo, también exige una arquitectura capaz de manejar secuencias largas, diferencias de escala y relaciones entre sensores.

% Una alternativa clásica consiste en extraer características temporales o estadísticas, como media, varianza, potencia, coeficientes autorregresivos o amplitud de potenciales relacionados con eventos. Estas representaciones reducen la dimensionalidad, pero dependen de decisiones manuales y pueden perder información. En habla imaginada se han utilizado también características de conectividad, patrones espaciales comunes y matrices de covarianza \cite{panachakel2021covert,alzahrani2024systematic}.

% Las representaciones tiempo-frecuencia describen cómo cambia el contenido espectral a lo largo de la señal. Para obtenerlas se emplean transformadas de Fourier de tiempo corto, bancos de filtros o transformadas wavelet. El resultado puede visualizarse como un espectrograma o escalograma. Estas imágenes permiten aplicar redes convolucionales desarrolladas originalmente para visión. No obstante, convertir una señal multicanal en una imagen obliga a decidir cómo se colocan los sensores y cómo se combinan sus dimensiones. Una representación inadecuada puede mezclar canales sin respetar su posición real.

% La dimensión espacial puede modelarse de varias formas. En EEG, la disposición de los electrodos sobre la cabeza se puede proyectar a un mapa bidimensional. En MEG, se dispone de la posición y orientación tridimensional de los sensores, además de su tipo. También pueden utilizarse matrices de covarianza o correlación para representar la relación entre canales. Otra posibilidad es reconstruir las fuentes corticales y trabajar en espacio cerebral, aunque este proceso requiere modelos anatómicos y añade una etapa de estimación.

% Los modelos recientes utilizan mecanismos de atención espacial. En lugar de asignar un peso fijo a cada canal, la red aprende qué sensores son más informativos para cada muestra. Este enfoque también permite proyectar conjuntos con diferente número de canales a una dimensión latente común. En decodificación de palabras con MEG, se ha utilizado atención espacial basada en las coordenadas de los sensores antes de aplicar convoluciones temporales \cite{dascoli2024words,ozdogan2025libribrain}.

% Otra estrategia consiste en añadir \textit{embeddings} que describen cada sensor. La posición, la orientación y el tipo de sensor se transforman en vectores y se suman a la representación de la señal. De esta forma, el modelo no recibe únicamente un índice de canal, sino información sobre el lugar en el que se realizó la medida. Esta idea es especialmente útil para combinar sistemas MEG heterogéneos y puede adaptarse a EEG utilizando las coordenadas de los electrodos.

% La representación temporal también depende de la longitud del contexto. Los métodos basados en palabras aisladas utilizan ventanas de unos pocos segundos y predicen cada ejemplo de forma independiente. Los modelos contextuales reciben secuencias de ventanas y pueden relacionar la respuesta de una palabra con las anteriores y posteriores. MEG-XL amplía esta idea durante el preentrenamiento, utilizando muestras continuas de dos minutos y medio. Sus resultados muestran que las representaciones aprendidas con contextos más largos se transfieren mejor a la clasificación de palabras \cite{jayalath2026megxl}.

% El contexto largo no implica conservar únicamente frecuencias lentas. Un filtro paso alto puede eliminar deriva de muy baja frecuencia y, aun así, el modelo puede aprender dependencias estadísticas entre eventos separados por varios segundos o minutos. La frecuencia de la señal indica la velocidad de sus oscilaciones, mientras que el contexto determina durante cuánto tiempo puede relacionar el modelo diferentes patrones. Son conceptos distintos y ambos influyen en la información disponible.

\section{Aprendizaje profundo para señales cerebrales}

% El aprendizaje profundo permite aprender representaciones directamente a partir de señales de alta dimensión. Frente a los métodos tradicionales, que dependen de características diseñadas manualmente, una red neuronal puede optimizar conjuntamente la extracción de patrones y la tarea final. Esta capacidad es especialmente útil en EEG y MEG, donde las relaciones relevantes pueden aparecer de forma distribuida entre canales, frecuencias y momentos temporales.

% Sin embargo, utilizar modelos más grandes no garantiza un mejor resultado. Las señales cerebrales presentan pocos ejemplos etiquetados, una variabilidad elevada y una relación señal-ruido reducida. Una arquitectura con demasiados parámetros puede memorizar sujetos, sesiones o estímulos. Por este motivo, el diseño del modelo debe acompañarse de regularización, particiones rigurosas, múltiples semillas y comparaciones frente a modelos sencillos.

\subsection{Redes neuronales y Transformers}

% Una red neuronal transforma una entrada mediante una sucesión de capas parametrizadas. Cada capa combina los valores de la anterior y aplica una función no lineal. Durante el entrenamiento, los parámetros se ajustan para reducir una función de pérdida. En clasificación, esta pérdida mide la diferencia entre las predicciones del modelo y las etiquetas reales. En aprendizaje de representaciones, puede medir la capacidad para reconstruir una señal o acercar ejemplos relacionados.

% Las redes totalmente conectadas pueden utilizarse sobre vectores de características, pero no aprovechan de forma explícita la estructura temporal. Las redes convolucionales, conocidas como CNN, aplican filtros locales que se desplazan por la señal. En EEG y MEG, una convolución temporal puede detectar patrones cortos que aparecen en distintos momentos. Las convoluciones dilatadas aumentan el campo receptivo sin necesitar un número excesivo de capas y se han utilizado como codificadores de señal en modelos de decodificación del habla \cite{defossez2023speech,dascoli2024words}.

% Las CNN también pueden operar sobre representaciones tiempo-frecuencia o mapas espaciales. Esta flexibilidad ha favorecido su uso en clasificación de habla imaginada. Su principal limitación es que las dependencias muy largas solo se modelan después de acumular numerosas capas o aumentar considerablemente el campo receptivo. Además, cuando se utilizan imágenes construidas artificialmente, el resultado depende de que la organización de los canales tenga un significado espacial adecuado.

% Las redes recurrentes procesan las muestras de forma secuencial y mantienen un estado interno. Arquitecturas como LSTM y GRU fueron diseñadas para conservar información a largo plazo mejor que las redes recurrentes simples. Se han aplicado a EEG, aunque su entrenamiento secuencial dificulta el paralelismo y puede resultar costoso con grabaciones largas.

% Los Transformers sustituyen la recurrencia por mecanismos de atención \cite{vaswani2017attention}. Para una secuencia de vectores, cada elemento se transforma en una consulta, una clave y un valor. La atención calcula la similitud entre consultas y claves y utiliza estos pesos para combinar los valores:
% \[
%     \operatorname{Attention}(\mathbf{Q},\mathbf{K},\mathbf{V}) =
%     \operatorname{softmax}\left(
%     \frac{\mathbf{Q}\mathbf{K}^{\top}}{\sqrt{d}}
%     \right)\mathbf{V}.
% \]
% Gracias a este mecanismo, una posición puede acceder directamente a cualquier otra posición de la secuencia. En decodificación cerebral, esto permite relacionar segmentos asociados a distintas palabras y seleccionar los momentos o sensores más informativos.

% El modelo necesita también información sobre el orden. Los \textit{embeddings} posicionales, ya sean aprendidos, sinusoidales o rotatorios, indican la posición temporal de cada elemento. Sin ellos, el mecanismo de atención trataría la secuencia como un conjunto sin orden. En clasificación contextual de palabras se han utilizado posiciones rotatorias, ya que permiten representar distancias relativas entre palabras y extenderse a secuencias de distinta longitud.

% La atención completa tiene un coste cuadrático respecto al número de elementos. En una señal cerebral, el número de elementos puede ser muy elevado porque depende del tiempo, los sensores y, en algunos modelos, los niveles de cuantización. Para reducir este coste, se puede factorizar la atención. La atención \textit{criss-cross} procesa por separado las relaciones temporales dentro de cada canal y las relaciones espaciales entre canales en cada instante. MEG-XL utiliza esta estrategia para modelar contextos de dos minutos y medio sin construir una matriz de atención completa sobre todos los tokens \cite{jayalath2026megxl}.

% Una arquitectura de este tipo conserva dos ejes distintos. La atención temporal aprende qué momentos de un mismo sensor están relacionados, mientras que la atención espacial aprende qué sensores deben combinarse en un instante determinado. Esta separación resulta natural para datos EEG y MEG, ya que evita comprimir prematuramente toda la estructura multicanal en un único vector.

\subsection{Aprendizaje autosupervisado y transferencia de aprendizaje}

% El aprendizaje supervisado necesita pares de señal y etiqueta. En decodificación de palabras, obtener estas etiquetas exige conocer el instante exacto de cada palabra y disponer de suficientes repeticiones. Los conjuntos de datos cerebrales suelen ser mucho más pequeños que los utilizados en visión o procesamiento del lenguaje. Además, registrar muchas horas de una misma persona es costoso y, en aplicaciones clínicas, puede ser inviable.

% El aprendizaje autosupervisado intenta aprovechar datos sin etiquetas manuales. La propia señal proporciona el objetivo de entrenamiento. Un ejemplo consiste en ocultar una parte de la entrada y pedir al modelo que la reconstruya a partir del contexto. Esta estrategia se utiliza en modelos de lenguaje enmascarado, audio y visión \cite{devlin2019bert,baevski2020wav2vec}. En señales neuronales, el modelo puede aprender regularidades temporales, relaciones entre canales y propiedades del dispositivo antes de recibir etiquetas lingüísticas.

% MEG-XL aplica predicción enmascarada sobre tokens discretos. Durante el preentrenamiento, se ocultan bloques temporales completos para todos los sensores y el Transformer debe predecir los códigos originales a partir de los bloques visibles \cite{jayalath2026megxl}. El enmascaramiento por bloques evita que la tarea se resuelva interpolando únicamente muestras contiguas. Además, al utilizar ventanas largas, el modelo debe aprender cuándo es útil atender a información cercana y cuándo debe integrar patrones más alejados.

% Después del preentrenamiento se realiza una adaptación o \textit{fine-tuning}. El modelo recibe ejemplos de la tarea final y sus parámetros se ajustan para predecir palabras. Si las representaciones aprendidas son útiles, el modelo necesita menos datos etiquetados que uno inicializado aleatoriamente. Otra forma de evaluación es el \textit{linear probing}, donde el modelo preentrenado se mantiene congelado y solo se entrena una capa lineal. Este procedimiento mide la calidad de las representaciones sin permitir que toda la red se adapte a la nueva tarea.

% La transferencia de aprendizaje puede realizarse entre tareas, sujetos, conjuntos de datos o modalidades. La transferencia entre sujetos pretende que patrones aprendidos con varias personas ayuden a un nuevo participante. La transferencia entre conjuntos introduce además diferencias de hardware, protocolo y estímulos. La transferencia de MEG a EEG es todavía más exigente, porque cambia la propia magnitud física registrada.

% Para que la transferencia sea efectiva, es necesario distinguir entre componentes generales y específicos. Las primeras capas pueden aprender patrones temporales básicos, mientras que las capas espaciales pueden depender fuertemente del sistema de sensores. Una estrategia consiste en reutilizar el Transformer temporal y adaptar la proyección de canales o los \textit{embeddings} de sensores. Otra posibilidad es ajustar toda la red con una tasa de aprendizaje pequeña, manteniendo una tasa mayor para la nueva cabeza de clasificación.

% También debe considerarse la transferencia negativa. Un modelo preentrenado puede rendir peor si las representaciones aprendidas no son adecuadas para el nuevo dominio. Por ejemplo, una frecuencia de muestreo, un filtrado o un tokenizer entrenado con una distribución diferente pueden descartar información útil. Por ello, la transferencia debe compararse con un modelo entrenado desde cero utilizando la misma arquitectura y los mismos datos. Solo así puede determinarse si la mejora procede realmente del preentrenamiento.

\subsection{Tokenización de señales neuronales}

% Los Transformers fueron desarrollados originalmente para secuencias de tokens lingüísticos. En texto, un token es un elemento discreto que identifica una palabra, una subpalabra o un símbolo. Las señales EEG y MEG, en cambio, son continuas. Para aplicar objetivos similares al modelado del lenguaje, se puede transformar la señal en una secuencia de representaciones discretas mediante un tokenizer neuronal.

% Un tokenizer aprende un codificador que comprime un fragmento de señal en un vector latente. A continuación, este vector se sustituye por la entrada más próxima de un diccionario o \textit{codebook}. El índice seleccionado actúa como token discreto. Un decodificador intenta reconstruir la señal original a partir de los tokens. Durante el entrenamiento del tokenizer se busca que la reconstrucción conserve la información relevante utilizando una secuencia más corta.

% La cuantización vectorial residual, conocida como RVQ, utiliza varios diccionarios de forma jerárquica. El primer nivel aproxima la señal y los siguientes codifican los residuos que quedan después de cada aproximación. Si se utilizan \(Q\) niveles, el fragmento queda representado por una secuencia de códigos
% \[
%     \mathbf{z} =
%     \left(z^{(1)},z^{(2)},\ldots,z^{(Q)}\right).
% \]
% Esta descomposición aumenta la capacidad de representación sin necesitar un único vocabulario de tamaño exponencial.

% BioCodec es un tokenizer desarrollado para bioseñales y entrenado principalmente con EEG \cite{avramidis2025biocodec}. En MEG-XL se aplica de forma independiente a cada canal y utiliza seis niveles RVQ, un vocabulario de 256 códigos por nivel y un factor de compresión temporal de doce \cite{jayalath2026megxl}. De esta forma, cada sensor produce varios tokens residuales por instante comprimido.

% La tokenización cumple dos funciones. En primer lugar, reduce la longitud temporal y permite procesar contextos mayores con el mismo presupuesto de memoria. En segundo lugar, proporciona un objetivo discreto para la predicción enmascarada. En lugar de reconstruir directamente cada valor continuo, el modelo predice una distribución de probabilidad sobre los códigos del vocabulario.

% Sin embargo, esta compresión también puede eliminar información. Un tokenizer entrenado con EEG a una frecuencia concreta no tiene por qué representar de forma óptima MEG, EEG remuestreado a otra frecuencia o bandas espectrales distintas. Por este motivo, la calidad de reconstrucción debe evaluarse en el dominio objetivo. Una reconstrucción visualmente parecida no garantiza que se conserven las características necesarias para clasificar palabras.

% También existen tokenizers que comprimen simultáneamente tiempo y canales. Esta estrategia produce secuencias más cortas, pero mezcla información espacial desde el principio. Si el objetivo es transferir entre sistemas con disposiciones de sensores diferentes, esta compresión puede dificultar la adaptación. La tokenización independiente por canal conserva la separación espacial, aunque obliga al Transformer posterior a modelar un número mayor de elementos.

% En el contexto de EEG, la selección del tokenizer está relacionada con la frecuencia de muestreo y el filtrado. Si la señal contiene frecuencias que el tokenizer nunca observó durante su entrenamiento, la cuantización puede ser poco precisa. Del mismo modo, reducir demasiado la frecuencia antes de tokenizar puede simplificar la señal a costa de eliminar componentes útiles. La comparación entre tokenizers, frecuencias y estrategias de inicialización constituye, por tanto, una parte importante de la adaptación de modelos MEG a EEG.

\section{Decodificación contextual de palabras}

% La clasificación de palabras intenta identificar qué palabra está asociada a un segmento de actividad cerebral. En el caso más sencillo, cada ventana se procesa de forma independiente y el modelo devuelve una distribución sobre un vocabulario fijo. En la decodificación contextual, el modelo recibe varias ventanas consecutivas y produce una predicción para cada una, teniendo en cuenta las relaciones dentro de la secuencia.

% Este enfoque se sitúa entre la clasificación aislada y la reconstrucción completa de texto. Sigue trabajando con un vocabulario cerrado y con eventos previamente alineados, pero introduce información lingüística y neural de las palabras vecinas. Esto permite estudiar si el contexto mejora la representación antes de añadir mecanismos más complejos de búsqueda, reordenación o generación de palabras fuera del vocabulario.

\subsection{Representaciones lingüísticas}

% Una clasificación convencional representa cada clase mediante un identificador independiente. En este caso, las palabras \textit{perro}, \textit{gato} y \textit{mesa} estarían igual de separadas entre sí, aunque las dos primeras compartan propiedades semánticas. Este planteamiento obliga al modelo a aprender cada palabra como una categoría aislada.

% Los \textit{embeddings} lingüísticos representan palabras mediante vectores continuos. Palabras utilizadas en contextos parecidos tienden a ocupar posiciones próximas en el espacio. Los modelos de lenguaje modernos generan representaciones contextuales, de modo que el vector asociado a una palabra también depende de la frase en la que aparece. Esta propiedad resulta útil para la decodificación cerebral, ya que la respuesta neuronal a una palabra tampoco es completamente independiente de su contexto.

% Los enfoques recientes utilizan representaciones extraídas de modelos T5 \cite{raffel2020t5}. Para cada palabra del vocabulario se obtiene un vector de una capa intermedia del modelo de lenguaje. Si una palabra se divide en varios tokens, sus representaciones pueden promediarse para construir un único objetivo a nivel de palabra. El decodificador cerebral no predice directamente una etiqueta, sino un vector que debe aproximarse al \textit{embedding} lingüístico correspondiente \cite{dascoli2024words,ozdogan2025libribrain}.

% Durante la inferencia, el vector neuronal predicho se compara con los vectores de todas las palabras del conjunto de recuperación. La palabra elegida es la que presenta mayor similitud coseno:
% \[
%     \hat{y} =
%     \arg\max_{y \in \mathcal{V}}
%     \frac{\hat{\mathbf{w}}^{\top}\mathbf{w}_{y}}
%     {\lVert\hat{\mathbf{w}}\rVert
%      \lVert\mathbf{w}_{y}\rVert},
% \]
% donde \(\hat{\mathbf{w}}\) es la representación predicha a partir de la señal cerebral, \(\mathbf{w}_{y}\) es la representación de la palabra \(y\) y \(\mathcal{V}\) es el vocabulario.

% Este planteamiento permite aprovechar la estructura semántica aprendida por el modelo de lenguaje. Una predicción incorrecta puede quedar cerca de palabras relacionadas en lugar de ser completamente arbitraria. Además, la misma arquitectura puede adaptarse a vocabularios diferentes cambiando el conjunto de vectores de recuperación. No obstante, el modelo continúa limitado por las palabras incluidas en dicho conjunto.

% El tamaño del vocabulario introduce un compromiso. Un vocabulario pequeño contiene menos opciones y facilita la clasificación, pero cubre una parte menor del texto. Un vocabulario grande se aproxima mejor al lenguaje natural, aunque reduce la precisión y necesita más ejemplos por clase. LibriBrain evalúa la clasificación sobre las 250 palabras más frecuentes, que cubren una proporción importante del contenido del corpus \cite{ozdogan2025libribrain}. MEG-XL utiliza también conjuntos más reducidos en algunos experimentos para analizar el rendimiento en escenarios con pocos datos \cite{jayalath2026megxl}.

% La frecuencia de las palabras es otro factor importante. Las palabras funcionales, como artículos y preposiciones, aparecen muchas veces y suelen clasificarse mejor que palabras de contenido menos frecuentes. Por ello, una precisión global puede estar dominada por unas pocas clases frecuentes. Las métricas equilibradas y los análisis por frecuencia ayudan a comprobar si el modelo aprende todo el vocabulario o únicamente las palabras más repetidas.

\subsection{Aprendizaje contrastivo}

% El aprendizaje contrastivo busca acercar representaciones correspondientes y separar representaciones no relacionadas. En la clasificación de palabras, cada segmento cerebral y su palabra real forman un par positivo. Las demás palabras del lote pueden utilizarse como ejemplos negativos. El modelo aprende un espacio común en el que la representación neuronal se aproxima a su objetivo lingüístico.

% Si \(\mathbf{r}_i\) es la representación cerebral del ejemplo \(i\) y \(\mathbf{w}_j\) es el \textit{embedding} de la palabra \(j\), se puede construir una matriz de similitudes
% \[
%     s_{ij} =
%     \frac{\mathbf{r}_i^{\top}\mathbf{w}_j}
%     {\lVert\mathbf{r}_i\rVert\lVert\mathbf{w}_j\rVert}.
% \]
% Los elementos de la diagonal corresponden normalmente a pares positivos y el resto actúa como negativos. La función de pérdida aumenta la similitud de los positivos y reduce la de los negativos.

% Una dificultad aparece cuando una misma palabra se repite dentro del lote. Dos segmentos asociados a la palabra \textit{the}, por ejemplo, no deberían considerarse negativos entre sí. Las variantes de SigLIP utilizadas en decodificación de palabras introducen máscaras para evitar que las repeticiones se penalicen incorrectamente \cite{zhai2023siglip,dascoli2024words}. MEG-XL emplea D-SigLIP durante el ajuste fino para predecir representaciones T5 \cite{jayalath2026megxl}.

% El aprendizaje contrastivo presenta varias ventajas. No obliga a utilizar una capa de clasificación con un tamaño fijo, incorpora relaciones entre palabras y permite reutilizar representaciones lingüísticas preentrenadas. También favorece el aprendizaje de una geometría común entre dos dominios muy diferentes: la señal cerebral y el lenguaje.

% Sin embargo, la elección de negativos afecta considerablemente al entrenamiento. Si los negativos son demasiado fáciles, el modelo puede aprender diferencias generales sin distinguir palabras parecidas. Si contienen falsos negativos, se penalizan relaciones válidas. El tamaño del lote, la temperatura de la similitud y el equilibrio de clases también influyen en el resultado.

% El contexto se incorpora antes de aplicar la pérdida. Cada ventana cerebral se codifica localmente y las representaciones de varias palabras consecutivas se introducen en un Transformer bidireccional. La salida correspondiente a cada posición contiene información de toda la secuencia. Después, cada una se compara con el \textit{embedding} lingüístico de su palabra. De este modo, el modelo aprende simultáneamente relaciones neuronales entre ventanas y relaciones semánticas entre palabras.

\subsection{Evaluación de la clasificación de palabras}

% La evaluación de un clasificador de palabras debe reflejar la dificultad del vocabulario y su distribución. La precisión Top-1 considera correcta una muestra únicamente cuando la palabra real ocupa la primera posición. La precisión Top-\(k\) considera correcta la muestra si la palabra aparece entre las \(k\) predicciones con mayor puntuación. En señales cerebrales, donde existe una incertidumbre elevada, Top-5 o Top-10 permiten medir si el modelo recupera un conjunto reducido de candidatos plausibles.

% Si el vocabulario contiene \(V\) palabras y todas las clases tienen la misma probabilidad, el nivel de azar para Top-\(k\) es aproximadamente
% \[
%     \operatorname{Azar}_{\text{Top-}k} = \frac{k}{V}.
% \]
% Por ejemplo, con 250 palabras, la probabilidad aleatoria de incluir la palabra correcta entre diez candidatos es \(10/250=0.04\). Esta referencia debe acompañar siempre a los resultados, aunque una línea base aleatoria basada en la distribución real puede ser más adecuada cuando las clases están desbalanceadas.

% La precisión equilibrada calcula primero el rendimiento de cada clase y después obtiene su media. Para Top-\(k\), puede expresarse como
% \[
%     \operatorname{BA}_{\text{Top-}k} =
%     \frac{1}{V}
%     \sum_{v=1}^{V}
%     \frac{\text{aciertos Top-}k\text{ de la clase }v}
%     {\text{número de ejemplos de la clase }v}.
% \]
% Esta métrica evita que las palabras más frecuentes dominen el resultado. Por este motivo, la precisión Top-10 equilibrada se ha convertido en una métrica habitual en clasificación contextual de palabras con MEG \cite{ozdogan2025libribrain,jayalath2026megxl}.

% Los resultados deben calcularse sobre datos no utilizados para ajustar el modelo ni seleccionar hiperparámetros. En estudios con varios sujetos es importante aclarar si la evaluación es dependiente o independiente del sujeto. En el primer caso, el modelo ha observado datos de esa persona durante el entrenamiento. En el segundo, se evalúa su capacidad para generalizar a participantes completamente nuevos. Esta segunda condición es más exigente y se aproxima mejor a una aplicación clínica.

% También deben evitarse fugas entre estímulos. Si la misma narración aparece en entrenamiento y prueba para sujetos distintos, el modelo puede reconocer características del audio, la duración o el orden de las palabras. Una partición rigurosa mantiene el mismo estímulo en una única división. Del mismo modo, las ventanas solapadas de una grabación no deben repartirse entre entrenamiento y prueba.

% Una evaluación fiable incluye comparaciones con varias líneas base. La primera es el nivel de azar. La segunda puede consistir en sustituir la señal cerebral por ruido aleatorio o permutar las correspondencias entre señal y palabra. Estos controles permiten comprobar si el modelo utiliza realmente información neuronal. También es conveniente comparar la arquitectura preentrenada con la misma arquitectura inicializada aleatoriamente. De esta forma se separa el efecto del preentrenamiento del efecto de la capacidad del modelo.

% La variabilidad entre ejecuciones debe reflejarse mediante varias semillas y medidas de dispersión, como desviación típica o error estándar. Para determinar si una mejora supera el azar se pueden utilizar pruebas de permutación. Cuando se comparan dos modelos sobre los mismos ejemplos, las pruebas deben tener en cuenta que las predicciones no son independientes.

% Además del resultado agregado, es útil analizar el rendimiento por frecuencia, categoría gramatical, sujeto y conjunto de datos. Estos análisis pueden revelar que la mejora se concentra en palabras frecuentes o en unos pocos participantes. También permiten estudiar si el contexto beneficia por igual a sustantivos, verbos y palabras funcionales.

% Por último, las métricas cambian cuando se pasa de clasificación de palabras a reconstrucción de secuencias. En sistemas \textit{brain-to-text} se utilizan medidas como tasa de error de palabra, tasa de error de carácter, BLEU, ROUGE o similitud semántica. Estas métricas evalúan la secuencia completa y pueden incorporar un modelo de lenguaje para reordenar o completar candidatos. Sin embargo, un buen resultado de secuencia no demuestra por sí solo que la señal cerebral sea informativa, ya que el modelo lingüístico puede corregir gran parte de la salida. Por ello, incluso en sistemas completos, siguen siendo necesarios controles con ruido, permutaciones y evaluación directa de la clasificación neuronal \cite{jayalath2025unlocking}.

% En conjunto, la clasificación contextual de palabras proporciona un marco controlado para estudiar la transferencia de representaciones entre MEG y EEG. Permite mantener una tarea, un espacio lingüístico y unas métricas comunes mientras se modifican la modalidad, el preprocesamiento, la frecuencia de muestreo, el tokenizer o la estrategia de inicialización. Esta separación facilita identificar qué componentes contribuyen realmente a la adaptación y qué limitaciones proceden de la calidad y cantidad de los datos.



\chapter{Estado del arte}

\section{Decodificación no invasiva del lenguaje}
\subsection{Decodificación mediante EEG}
\subsection{Decodificación mediante MEG}

\section{Datasets de lenguaje con EEG y MEG}
\subsection{Datasets de escucha}
\subsection{Datasets de lectura}

\section{Modelos preentrenados de señales cerebrales}
\subsection{Modelos para EEG}
\subsection{MEG-XL}

\section{Limitaciones de los trabajos existentes}


\chapter{Propuesta: adaptación de MEG-XL a EEG}

\section{Descripción general de la propuesta}

\section{Preparación de los datasets}
\subsection{Datasets de lectura}
\subsection{Datasets de escucha}
\subsection{Representación unificada de los datos}

\section{Preprocesamiento del EEG}
\subsection{Filtrado y remuestreo}
\subsection{Normalización y segmentación}
\subsection{Alineamiento con las palabras}

\section{Adaptación de la arquitectura MEG-XL}
\subsection{Adaptación de los canales EEG}
\subsection{Tokenización de la señal}
\subsection{Modelado temporal y espacial}

\section{Entrenamiento del modelo}
\subsection{Preentrenamiento autosupervisado}
\subsection{Clasificación contextual de palabras}
\subsection{Inicialización desde cero y desde MEG-XL}

\section{Implementación del sistema}
\subsection{Organización del repositorio}
\subsection{Ejecución y monitorización de experimentos}


\chapter{Evaluación experimental}

\section{Configuración de los experimentos}
\subsection{Particiones de los datos}
\subsection{Hiperparámetros y métricas}

\section{Influencia de la frecuencia de muestreo}

\section{Influencia de las bandas de frecuencia}

\section{Entrenamiento desde cero y transferencia desde MEG-XL}


\section{Comparación de tokenizadores}

\section{Comparación de tareas}
\subsection{Lectura}
\subsection{Escucha}
\subsection{Lectura y escucha combinadas}

\section{Discusión de los resultados}
\subsection{Principales conclusiones experimentales}
\subsection{Limitaciones}

%%%%%%%%%%%%%%%%%%%%%%%%%%%%%%%%%%%%%%%%%%%%%%%%%%%%%%%%%%%%%%%%%%%%%%%%%%%%%%%
%                                 CONCLUSIONS                                 %
%%%%%%%%%%%%%%%%%%%%%%%%%%%%%%%%%%%%%%%%%%%%%%%%%%%%%%%%%%%%%%%%%%%%%%%%%%%%%%%

\chapter{Conclusiones y trabajo futuro}

\section{Conclusiones}

\section{Cumplimiento de los objetivos}

\section{Líneas de trabajo futuro}

%%%%%%%%%%%%%%%%%%%%%%%%%%%%%%%%%%%%%%%%%%%%%%%%%%%%%%%%%%%%%%%%%%%%%%%%%%%%%%%
%                                BIBLIOGRAFIA                                 %
%%%%%%%%%%%%%%%%%%%%%%%%%%%%%%%%%%%%%%%%%%%%%%%%%%%%%%%%%%%%%%%%%%%%%%%%%%%%%%%

\begin{thebibliography}{10}

%%%%%%%%%%%%%%%%%%%%%%%%%%%%%%%%%%%%%%%%%%%%%%%%%%%%%%%%%%%%%%%%%%%%%%%%%%%%%%%
% MODEL D'ARTICLE                                                             %
%%%%%%%%%%%%%%%%%%%%%%%%%%%%%%%%%%%%%%%%%%%%%%%%%%%%%%%%%%%%%%%%%%%%%%%%%%%%%%%
\bibitem{light}
   Jennifer~S. Light.
   \newblock When computers were women.
   \newblock \textit{Technology and Culture}, 40:3:455--483, juliol, 1999.

%%%%%%%%%%%%%%%%%%%%%%%%%%%%%%%%%%%%%%%%%%%%%%%%%%%%%%%%%%%%%%%%%%%%%%%%%%%%%%%
% MODEL DE LLIBRE                                                             %
%%%%%%%%%%%%%%%%%%%%%%%%%%%%%%%%%%%%%%%%%%%%%%%%%%%%%%%%%%%%%%%%%%%%%%%%%%%%%%%
\bibitem{ifrah}
   Georges Ifrah.
   \newblock \textit{Historia universal de las cifras}.
   \newblock Espasa Calpe, S.A., Madrid, sisena edició, 2008.

%%%%%%%%%%%%%%%%%%%%%%%%%%%%%%%%%%%%%%%%%%%%%%%%%%%%%%%%%%%%%%%%%%%%%%%%%%%%%%%
% MODEL D'URL                                                                 %
%%%%%%%%%%%%%%%%%%%%%%%%%%%%%%%%%%%%%%%%%%%%%%%%%%%%%%%%%%%%%%%%%%%%%%%%%%%%%%%
\bibitem{WAR}
   Comunicat de premsa del Departament de la Guerra, 
   emés el 16 de febrer de 1946. 
   \newblock Consultat a 
   \url{http://americanhistory.si.edu/comphist/pr1.pdf}.

\end{thebibliography}
\cleardoublepage

%%%%%%%%%%%%%%%%%%%%%%%%%%%%%%%%%%%%%%%%%%%%%%%%%%%%%%%%%%%%%%%%%%%%%%%%%%%%%%%
%                           APÈNDIXS  (Si n'hi ha!)                           %
%%%%%%%%%%%%%%%%%%%%%%%%%%%%%%%%%%%%%%%%%%%%%%%%%%%%%%%%%%%%%%%%%%%%%%%%%%%%%%%

\APPENDIX

%%%%%%%%%%%%%%%%%%%%%%%%%%%%%%%%%%%%%%%%%%%%%%%%%%%%%%%%%%%%%%%%%%%%%%%%%%%%%%%
%                         LA CONFIGURACIO DEL SISTEMA                         %
%%%%%%%%%%%%%%%%%%%%%%%%%%%%%%%%%%%%%%%%%%%%%%%%%%%%%%%%%%%%%%%%%%%%%%%%%%%%%%%

\chapter{Configuració del sistema}

????? ????????????? ????????????? ????????????? ????????????? ?????????????

\section{Fase d'inicialització}

????? ????????????? ????????????? ????????????? ????????????? ?????????????

\section{Identificació de dispositius}

????? ????????????? ????????????? ????????????? ????????????? ?????????????

%%%%%%%%%%%%%%%%%%%%%%%%%%%%%%%%%%%%%%%%%%%%%%%%%%%%%%%%%%%%%%%%%%%%%%%%%%%%%%%
%                               ALTRES  APÈNDIXS                              %
%%%%%%%%%%%%%%%%%%%%%%%%%%%%%%%%%%%%%%%%%%%%%%%%%%%%%%%%%%%%%%%%%%%%%%%%%%%%%%%


\chapter{??? ???????????? ????}

????? ????????????? ????????????? ????????????? ????????????? ????????????? 



%%%%%%%%%%%%%%%%%%%%%%%%%%%%%%%%%%%%%%%%%%%%%%%%%%%%%%%%%%%%%%%%%%%%%%%%%%%%%%%
%                              FI DEL DOCUMENT                                %
%%%%%%%%%%%%%%%%%%%%%%%%%%%%%%%%%%%%%%%%%%%%%%%%%%%%%%%%%%%%%%%%%%%%%%%%%%%%%%%

\end{document}
